\documentclass[preview]{standalone}

\usepackage{amsmath}
\usepackage{amssymb}
\usepackage{stellar}
\usepackage{definitions}

\begin{document}

\id{field-extensions}
\genpage

\section{Field extensions}

\begin{snippetdefinition}{field-extension-definition}{Field extension}
    Let \(K\) and \(L\) be \field[fields]. A \emph{field extension of
    \(K\)} is a \field \(L\) equipped with an injective
    \ringhomomorphism
    \[
        \iota\colon K\fromto L.
    \]
    After identifying \(K\) with \(\iota(K)\subseteq L\), one writes
    \[
        L/K.
    \]
\end{snippetdefinition}

\begin{snippetdefinition}{field-extension-degree-definition}{Degree of a field extension}
    Let \(L/K\) be a \fieldextension. Multiplication in \(L\) makes \(L\)
    a \vectorspace over \(K\). The \emph{degree} of the extension is
    \[
        [L:K]=\lineardim_K L.
    \]
    The extension is \emph{finite} if \([L:K]<\infty\).
\end{snippetdefinition}

\section{Extensions from irreducible polynomials}

\begin{snippetproposition}{irreducible-polynomial-quotient-field-extension}{Quotient by an irreducible polynomial}
    Let \(K\) be a \field and let \(f\in K[x]\) be an irreducible
    \polynomial of degree \(n\geq1\). Then
    \[
        L=K[x]/(f)
    \]
    is a \fieldextension of \(K\) of degree \(n\). A \(K\)-basis is
    \[
        \mathcal{B}
        =
        \{1+(f),x+(f),\ldots,x^{n-1}+(f)\}.
    \]
\end{snippetproposition}

\begin{snippetproof}{irreducible-polynomial-quotient-field-extension-proof}{irreducible-polynomial-quotient-field-extension}{Quotient by an irreducible polynomial}
    Since \(K[x]\) is a \principalidealdomain, the irreducibility of
    \(f\) makes \((f)\) a \maximalideal. Hence \(L=K[x]/(f)\) is a
    \field.

    The map
    \[
        \iota\colon K\fromto L,
        \qquad
        \iota(\lambda)=\lambda+(f),
    \]
    is a \ringhomomorphism. If \(\iota(\lambda)=0_L\), then the constant
    polynomial \(\lambda\) belongs to \((f)\). Since
    \(\polynomialdeg f\geq1\), this forces \(\lambda=0_K\), so \(\iota\)
    is injective and \(L/K\) is a \fieldextension.

    Every element of \(L\) has the form \(g+(f)\) for some \(g\in K[x]\).
    Polynomial division gives
    \[
        g=qf+r,
        \qquad
        r=0\text{ or }\polynomialdeg r<n.
    \]
    Thus \(g+(f)=r+(f)\), and writing
    \[
        r=a_0+a_1x+\cdots+a_{n-1}x^{n-1}
    \]
    shows that \(\mathcal{B}\) spans \(L\).

    If
    \[
        a_0(1+(f))+\cdots+a_{n-1}(x^{n-1}+(f))=0_L,
    \]
    then
    \[
        a_0+\cdots+a_{n-1}x^{n-1}\in(f).
    \]
    A nonzero member of \((f)\) has degree at least \(n\), so all
    coefficients \(a_i\) vanish. Therefore \(\mathcal{B}\) is linearly
    independent and \([L:K]=n\).
\end{snippetproof}

\section{Examples}

\begin{snippetexample}{gaussian-rational-field-extension-example}{The field of Gaussian rationals}
    The polynomial \(x^2+1\in\rationalnumbers[x]\) is irreducible because
    it has no rational root. Hence
    \[
        L=\rationalnumbers[x]/(x^2+1)
    \]
    is a degree-two extension of \(\rationalnumbers\), and every element
    has a unique representative \(a+bx\), with
    \(a,b\in\rationalnumbers\).

    In \(L\),
    \[
        (x+(x^2+1))^2=-1+(x^2+1),
    \]
    so \(x+(x^2+1)\) behaves like \(i\). More precisely, evaluation at
    \(i\) gives a surjective \ringhomomorphism
    \[
        \varphi\colon\rationalnumbers[x]\fromto\rationalnumbers(i),
        \qquad
        \varphi(g)=g(i).
    \]
    Dividing \(g\) by \(x^2+1\) shows that \(g(i)=0\)
    \ifandonlyif \(x^2+1\divides g\). Therefore
    \[
        \ringker\varphi=(x^2+1),
    \]
    and the first isomorphism theorem gives
    \[
        \rationalnumbers[x]/(x^2+1)
        \ringisomorphic
        \rationalnumbers(i).
    \]
\end{snippetexample}

\begin{snippetexample}{gaussian-rational-field-inverses-example}{Inverses in the Gaussian rational field}
    Let \(I=(x^2+1)\idealin\rationalnumbers[x]\). Every nonzero element of
    \(\rationalnumbers[x]/I\) has the form
    \[
        (a+bx)+I,
        \qquad
        a,b\in\rationalnumbers,
        \quad
        (a,b)\neq(0,0).
    \]
    Its inverse is
    \[
        \bigl((a+bx)+I\bigr)^{-1}
        =
        \frac{a-bx}{a^2+b^2}+I.
    \]
    Indeed,
    \begin{align*}
        \bigl((a+bx)+I\bigr)
        \left(\frac{a-bx}{a^2+b^2}+I\right)
        &=
        \frac{a^2-b^2x^2}{a^2+b^2}+I\\
        &=
        1+I,
    \end{align*}
    because \(x^2+I=-1+I\).
\end{snippetexample}

\begin{snippetexample}{irreducible-polynomial-quotient-inverse-example}{An inverse in a polynomial quotient}
    Let
    \[
        f=x^4+3x^3+x^2-2x+1\in\rationalnumbers[x].
    \]
    The polynomial \(f\) is irreducible, so
    \(L=\rationalnumbers[x]/(f)\) is a \field. To invert
    \(g+(f)\), where \(g=x^2+1\), apply the Euclidean algorithm:
    \[
        f=(x^2+3x)g+(1-5x)
    \]
    and
    \[
        g=
        \left(-\frac15x-\frac1{25}\right)(1-5x)+\frac{26}{25}.
    \]
    Back-substitution gives
    \[
        1
        =
        \left[
            \frac{25}{26}
            -(x^2+3x)\left(\frac5{26}x+\frac1{26}\right)
        \right]g
        +
        \left(\frac5{26}x+\frac1{26}\right)f.
    \]
    Consequently,
    \[
        (g+(f))^{-1}
        =
        \left[
            \frac{25}{26}
            -(x^2+3x)\left(\frac5{26}x+\frac1{26}\right)
        \right]+(f).
    \]
\end{snippetexample}

\end{document}
